\section{\acs{5G} Resource Lifecycle Management}
\label{chap:impl:resource_lifecycle}

As established in Section~\ref{chap:proposal:resource_lifecycle}, network slicing is the main mechanism that allows the testbed to support diverse vertical use cases and network applications, and their respective heterogeneous \acs{QoS} requirements. Delivering this requires the automation and orchestration of each phase of the \acs{NSI} lifecycle, a responsibility handled by the coordinated operation of the \acs{CSMF}, \acs{NSMF} and \acs{NSSMF} management functions. In this section, the components responsible for realizing each of these management functions are identified. However, in practice, not all of these functions have a direct counterpart within the implementation developed in this dissertation, being instead abstracted by other components already present in the \acs{ITAv} 5GAIner testbed~\cite{5GAIner}.

At the top of this management hierarchy, the \acs{OSL} \acs{OSS}~\cite{OpenSlice} is responsible for the \acs{OAM} functionality within the testbed, fulfilling the \acs{CSMF} role established in Section~\ref{chap:proposal:resource_lifecycle}. In this capacity, it exposes communication services to third-party clients through \acs{TMF} Open \acsp{API}, allowing them to request and provision \acs{5G} resources (network slices and \acsp{UE}) without requiring privileged knowledge of the underlying infrastructure. This abstraction is implemented by leveraging the \acf{CFS}, \acf{RFS} and Resource constructs, defined by \acs{TMF}, as mentioned in Section~\ref{sota:telcoinitiatives:tmf}. More concretely, each \acs{CFS} is exposed to customers as a configurable interface, through which clients can specify their desired service parameters (e.g. uplink throughput for a slice). Internally, the fields of the \acs{CFS}, expressing the user's requirements, are then mapped into a corresponding \acs{RFS} which in turn references the concrete Resource capturing the technical parameters required to provision the underlying network slice or \acs{UE}.

However, as previously discussed, \acs{OSL} is not an orchestrator, meaning the provisioning of the resources referenced by each \acs{RFS} is delegated to dedicated components. Therefore, \acs{OSL}, can delegate the provisioning of \acs{NFV} resources, to \acs{NFV} \acs{OSM}, via the SOL005 interface. Containerized resources, such as network applications, can also be automatically deployed to kubernetes through the CRIDGE component. In the case of network slices and \acsp{UE}, the provisioning of resources is delegated to a proprietary Slice Manager within \acs{ITAv}'s testbed.

This proprietary Slice Manager represents the primary mechanism for interacting with the \acs{5G} System, both the \acs{RAN} and the core, to provision network slices and \acsp{UE} in the 5GAIner infrastructure. This component abstracts the underlying complexity of interacting with the \acs{5G} System and \acs{NFV} infrastructure behind a unified \acs{API}, simplifying the management of \acs{5G} resources. Given the abstraction provided by this \acs{API}, the design and implementation of the \acs{NSMF} and \acs{NSSMF} services falls outside the scope of this dissertation, as their underlying functionality is already encapsulated within the Slice Manager. Consequently, \acs{OSL} must be capable of converting high-level customer service requirements, provided via a \acs{CFS}, into concrete \acs{API} calls to the Slice Manager, resulting in the provisioning of network slices and \acsp{UE}. This can be accomplished through the use of \acs{OSL} Generic Controllers~\cite{OpenSliceGenericServicesController}, which will be discussed in detail in the next section.

\subsection{Slicing Generic Controller}

A Generic Controller, inspired by the kubernetes operator pattern, is an \acs{OSL} component capable of managing resources of a specific type (of their resource specification)~\cite{OpenSliceGenericServicesController}. It is integrated natively into the OSL lifecycle, allowing custom code to be executed in response to the creation, update, and deletion of \acs{OSL} resources, by listening to a dedicated message queue for resource lifecycle events. This automated process allows the \acs{OSL} resources and services to be managed programmatically, ensuring they remain aligned with the state of the \acs{5G} deployment.

The Generic Controller developed in this dissertation manages two resource types: network slices and \acsp{UE}. For each resource type, a Resource Specification was defined, establishing the schema of the managed resource, aligned with the corresponding resource fields required by the Slice Manager \acs{API}. The controller runs as a standalone container, integrated with the \acs{OSL} deployment via its message queue. In this section, the design and implementation of the Generic Controller is presented. The lifecycle management of network slices and \acsp{UE} is covered first, detailing how the controller translates \acs{OSL} service events into Slice Manager \acs{API} calls for each resource type. The service composition mechanisms that coordinate the customer-facing and resource-facing service layers are then described. Finally, the integration with the testbed observability infrastructure is addressed.

\plan{Make a better introduction making a broad arrangement of the following subsubsections}

\subsubsection{Network Slice Lifecycle Management}

The Slicing Generic Controller manages the lifecycle of network slice resources in the testbed. Each slice resource is defined by a Resource Specification capturing the parameters required for provisioning the resource using the Slice Manager \acs{API}. The required payload is provided in Appendix~\ref{appendix:slicemanager:payloads}, Code Block~\ref{code:slicemanager:payloads:netslice}. Its most relevant parameters include:

\begin{itemize}
    \item \texttt{id} \texttt{|} \texttt{name}: uniquely identify the slice resource in the \acs{5G} System;
    \item \texttt{administrative\_state} \texttt{|} \texttt{operational\_state}: indicate the slice status, regarding service (\texttt{LOCKED}, \texttt{UNLOCKED} or \texttt{SHUTTINGDOWN}) and operational state (\texttt{ENABLED} or \texttt{DISABLED}), respectively. The former controls whether the slice is available for use, while the latter its actual runtime condition and if it's capable of receiving traffic;
    \item \texttt{sst} \texttt{|} \texttt{ssd}: collectively form the \acs{S-NSSAI} identifier, classifying the slice type (e.g. \acs{eMBB}) and differentiating between slices of the same type, respectively;
    \item \texttt{dnn}: identifies the name of the target data network for the network slice;
    \item \texttt{dllatency} \texttt{|} \texttt{ullatency} \texttt{|} \texttt{dlguathptperue} \texttt{|} \texttt{ulguathptperue} \texttt{|} \texttt{dlmaxthpt\-perue} \texttt{|} \texttt{ulmaxthptperue} \texttt{|} \texttt{dlguathptperslice} \texttt{|} \texttt{ulguathptperslice} \texttt{|} \texttt{dl\-maxthptperslice} \texttt{|} \texttt{ulmaxthptperslice}: define the \acs{QoS} performance requirements of the slice, covering downlink and uplink latency, guaranteed and maximum throughput per \acs{UE}, and guaranteed and maximum throughput of the slice, respectively.
\end{itemize}

The fields of the Resource Specification, illustrated in Figure~\ref{} and defined in Code Block ~\ref{} (at Appendix~\ref{}), are aligned with the parameters accepted by the Slice Manager \acs{API}, allowing the Generic Controller to directly incorporate the resource state into the corresponding \acs{API} payload. Upon resource creation, the Generic Controller receives the  corresponding \texttt{CREATION} message and issues a \texttt{POST} request to the \texttt{/productOrder} endpoint of the Slice Manager \acs{API}, instantiating the network slice across the \acs{RAN} and \acs{CN}. Update events trigger a \texttt{PATCH} request to the \texttt{/productOrder/\{id\}/patch} endpoint of the Slice Manager \acs{API}, including the slice id and the modified resource parameters to update the slice configuration according to the new client requirements. Deletion of the resource triggers a \texttt{DELETE} request to the \texttt{/productOrder/\{id\}} endpoint of the Slice Manager \acs{API}, resulting in the deprovisioning of the slice and disassociation of any linked \acsp{UE}, operations that are once again abstracted by the Slice Manager component.

\subsubsection{\acl{UE} Lifecycle Management}

\plan{Brief transition: provisioning a slice is necessary but insufficient for connectivity — a UE must be associated with the slice for traffic to flow through it. The Generic Controller handles this association as part of the UE resource lifecycle.}

\subsubsection{Service Composition}

\subsubsection{Observability Integration}

The network slice component manages the lifecycle of network slice resources in the testbed. Each slice resource is defined by a Resource Specification capturing the service parameters required for provisioning, such as the slice type and associated \acs{QoS} targets. Upon resource creation, the Generic Controller translates these parameters into Slice Manager \acs{API} calls to instantiate the network slice across the \acs{RAN} and \acs{CN}, realising the \acs{NSI} Automated Configuration procedure discussed in Section~\ref{chap:proposal:resource_lifecycle}. Update events trigger the corresponding Automated Reconfiguration procedure, propagating modified service parameters to the Slice Manager. Deletion events initiate the decommissioning of the slice, instructing the Slice Manager to release the associated network resources.

The \acs{UE} component manages the lifecycle of \acs{UE} resources and their association with provisioned network slices. Each \acs{UE} resource is defined by a Resource Specification capturing the identity and configuration parameters of the device. Upon resource creation, the Generic Controller interacts with the Slice Manager \acs{API} to register and associate the \acs{UE} with the target slice, configuring the required session parameters as described in Section~\ref{chap:proposal:resource_lifecycle}. Modifying the slice association of an active \acs{UE} is not supported in the current implementation, as the Slice Manager \acs{API} does not expose a dedicated endpoint for this operation. Deletion of a \acs{UE} resource triggers a single Slice Manager \acs{API} call that disassociates the \acs{UE} from its slice and releases its network configuration.




The Generic Controller is scoped to the management of Slice and \acs{UE} Resources and their corresponding \acsp{RFS}, and does not directly interact with the customer-facing service layer. LCM Rules fulfil this bridging role: when a customer updates \acs{CFS} parameters, the corresponding LCM Rule propagates the changed fields to the underlying \acs{RFS}, which the Generic Controller then acts upon. Conversely, when the Generic Controller updates the state of an \acs{RFS} following a provisioning operation, a complementary LCM Rule propagates this state back to the \acs{CFS}, ensuring the customer-visible service reflects the actual state of the provisioned resource.

Beyond resource provisioning, the automated lifecycle management must also integrate monitoring mechanisms, as established in Section~\ref{chap:proposal:observability}. Without visibility into the runtime state of provisioned resources, it is not possible to reliably attribute observed network application behaviour to the application logic rather than to transient infrastructure conditions. Accordingly, per-\acs{UE} telemetry collection is activated within the Generic Controller as part of the provisioning lifecycle: upon creation of a \acs{UE} resource, the controller triggers the activation of the corresponding Prometheus exporters, scoping telemetry collection to the \acsp{UE} associated with the provisioned slice. The implementation of the full monitoring stack is discussed in detail in Section~\ref{chap:impl:observability}.

As part of this monitoring integration, a Grafana dashboard is automatically provisioned for each slice, aggregating throughput metrics for all \acsp{UE} belonging to that slice. These metrics are passively collected by Prometheus exporters deployed across the network, capturing per-\acs{IP} and per-\acs{MAC} source throughput from the traffic of all network \acsp{UE}, including those associated with the provisioned slice.

\plan{\begin{enumerate}
    \item Introduce the implemented generic controller to manage the lifecycle of network slices and UEs (and associate a network slice to an UE). Mention generic controller permits the registration of Resource Specifications, and definition of creation, update and deletion callbacks to manage the corresponding resoures. Highlight that is the responsibility of the operator (in this case us, the owners of the testbed) to create a resource and resource specification for each of the 5G resources. The generic controller runs as a separate container that will manage each of this resources. In order to expose this resource to third parties a Resource Facing Service (and specification) and Customer Facing Service (and specification) must be created and coordinated (exchange of information between them using LCM Rules)
    \item Detail the network slice management component of the Generic Controller. Start by defining the Resource (and its fields). Then, specify how the generic controller handles the creation, update, and deletion of network slices translating OSL service requirements (fields) into the corresponding Slice Manager API calls, realizing the SON Automated Configuration and Reconfiguration procedures discussed in Section~\ref{chap:proposal:resource_lifecycle}
    \item Detail the UE management component of the Generic Controller. Start by defining the Resource (and its fields). Then, specify how a UE can be created and associated to a given slice. Highlight how in an optimal setting it would be possible to change the slice an UE belongs to (but it was not possible due to the absence of such an endpoint in the Slice Manager API). Mention the delete procedure and how Slice Manager API abstracts the disassociation of the UE from the respective slice under one API call.
    \item Highlight that the Generic Controller is only responsible for the management of the Slice/UE Resource and the corresponding Resource Facing Service. Therefore it is necessary to use LCM RULES to propagate the requirements/fields of the CFS to the RFS and vice-versa.
    \item Reference section 3.4, reenforcing the need for a fully automated monitoring process. Mention the implemented monitoring process will be further explained at Section \ref{chap:impl:observability}. 
    \item Describe the integration of the monitoring stack (Prometheus exporters and Grafana) with the lifecycle management, detailing how per-UE telemetry collection is activated within the Generic Controllers as part of the provisioning lifecycle, supporting the attribution of observed network application behaviour to the application logic rather than to infrastructure conditions.
    \item Mention the automatic provisioning of a grafana dashboard for all slice UEs with per-IP and per-MAC source throughput metrics, passively collected by the deployed exporters across all network UEs (and therefore that slice-UEs as well).
    \item Mention the dashboard also contains active probing metrics for the slice UEs, which requires publishing the UEs IP addresses, to the prober container, also described in Section \ref{chap:impl:observability}. Reenforce that this will only happen if the client, definining its requirements through the Slice Customer Facing Service, actually sets the active\_probing field to true. 
    \item Highlight the need for OSL service changes to be actively propagated to the network, reflecting the current state of 5G resources, active slices and connected UEs, and for the outcome, whether success or failure, to be reported back to the user. When the customer changes the criteria for the slice then it must be reflected onto the network. When the slice and UE are terminated so does that need to happen in the network. If any of the operations doesn't go as expected the showcased OSL service must not be offsync compared to the state of the network.
\end{enumerate}}